\honest{The Cluster Structure of Thingspace}{Eliezer Yudkowsky}{2008}

Colours can be written as coordinates, blue, blue-green and red as 0:0:5, 0:3:2 and 5:0:0, and then you can see how much closer blue is to blue-green than to red. In the same way I picture a robin as a point in ``thingspace,'' a space with a dimension for everything that could be known about it: volume (more than a virus, less than an aircraft carrier), mass, DNA as millions of four-valued coordinates, shape, colour, and redundantly density too. Across robins, volume and mass line up along a rough line, which is easier to see among points than among birds. Nothing is lost in the translation: a robin balanced against 0.07 kilograms and a robin-point at mass +70 say the same thing. If you think such a space is extravagant, quantum physicists use an infinite-dimensional one in which a single point describes every particle in the universe. If I am unsure of a robin's mass and volume, the robin becomes a cloud of probability; if I am surer of its density, the cloud concentrates along a slanting line, since volume times density is mass.

Then I turn to ``radial categories,'' which is how ``cognitive psychologists describe the non-Aristotelian boundaries of words.'' The central mother conceives, bears and supports her child; genetic, surrogate and adoptive mothers are variants. Logic says that humans have ten fingers and Fred has nine, so Fred is not human; we actually conclude that Fred is a ``nine-fingered human.'' \nb{The term and the example are George Lakoff's, from \textsc{Women, Fire, and Dangerous Things} (1987), which the previous post cites.} The meaning of ``mother'' might be pictured as a glow in thingspace, brightest at that centre, dimmer where the egg donors are. Map the world's birds with distances matching perceived similarity, robins closer to robins than to pigeons, and all of them closer to each other than to penguins, and they form something like an astronomical cluster: robins, sparrows, canaries and pigeons packed at the centre, eagles and falcons nearby, penguins, chickens and ostriches farther out. Picture both at once, the central clusters glowing brightly, the outliers dimly, and Abraham Lincoln megaparsecs away, not glowing at all.

``As I see it,'' the glow was learned from the clusters: first the structure in the world, then the category. \nb{The birds were placed by ``perceived similarity in humans,'' so the clusters are human judgments and cannot show which came first. Eleanor Rosch, stating the same principle in 1978, limited it to ``the perceived world'' and said that existing categories may help decide which attributes people perceive.}

So a natural cluster may have no set of necessary and sufficient properties. \nb{This is Wittgenstein's ``family resemblance'' (1953).} A simple definition may meet exceptions, and that is fine; you may find one for every simple rule. The map is smaller than the territory, and a definition's job is to lead the listener to the cluster, not to describe every bird to the molecule. ``Feathered flying things'' leads you to the birds even though penguins do not fly, and objecting to it is ``nitpicking.''
